% Table S2 — Year-level Pearson correlations between yield and candidate predictors
% VALUES VERIFIED against year_level_correlations.csv (source of truth)
\begin{table}[H]
\caption{Year-level Pearson correlations between mean olive yield and candidate
climate predictors ($n=8$ years). $T_{\text{chill},t-1}$ ($r=-0.697$) is the
strongest single predictor; $T_{\text{CREC},t}$ was selected as the heat-side
predictor for its low collinearity with $T_{\text{chill},t-1}$ ($r=0.409$,
VIF\,$<2$) and its clean physiological decomposition. Threshold-based chill
counts (CP, CH$_{12}$) are weaker because they saturate at zero in the warm tail.\label{tab:year_corr}}
\centering
\small
\begin{tabular}{lcc}
\toprule
\textbf{Predictor} & \textbf{Pearson $r$ vs yield} & \textbf{Note} \\
\midrule
$T_{\text{chill},t-1}$ (May--Aug, $^\circ$C)     & $\mathbf{-0.697}$ & primary (chill side) \\
$T_{\max,\text{FLOR}}$ ($^\circ$C)               & $-0.492$ & flowering $T_{\max}$ \\
CH$_{12,t-1}$ (h)                                & $+0.404$ & threshold chill, h$<12\,^\circ$C \\
$T_{\text{mean},\text{FLOR}}$ ($^\circ$C)        & $-0.381$ & flowering window \\
$T_{\max,\text{CREC}}$ ($^\circ$C)               & $-0.375$ & fruit growth $T_{\max}$ \\
$T_{\text{CREC},t}$ (Jan--Mar, $^\circ$C)        & $\mathbf{-0.313}$ & primary (heat side) \\
CP$_{t-1}$ (Dynamic Model)                       & $+0.307$ & saturated at 0 in 5/8 years \\
ET$_{0,\text{CREC}}$ (mm)                        & $-0.178$ & FAO~56 reference ET, fruit growth \\
\midrule
\multicolumn{3}{l}{\textit{Cross-correlation $T_{\text{chill},t-1}\times T_{\text{CREC},t}=+0.409$ (moderate; VIF\,$<2$)}} \\
\bottomrule
\end{tabular}
\end{table}
